% Solutions to Chapter 7's exercises, from lectures/L07/appendix/e_solutions.md. Exercise 6 is a
% program: the program itself is left out, and the parts worked on paper, the register values and
% the cycle count, are kept. The one or two lines of assembly in exercises 1, 3 and 5 are kept,
% because there the lines are the answer, and exercise 5 asks for the corrected program.
%
% Exercise 10 a) says that the program sits in the template, as the exercise now does; see
% chapters/07/exercises.tex.

\answerchapter{c7:ch}

\answer{c7:ex:1}{Arbetsbänken}
\pt{a}32 register, \code{r0}--\code{r31}, med 8 bitar, en byte, vardera.

\pt{b}\code{r16}--\code{r31}.

\pt{c}\code{ldi} är 16 bitar lång. Fyra av bitarna är operationskoden och åtta är konstanten, så
det blir bara \textbf{fyra bitar} kvar till registret. Fyra bitar kan räkna från 0 till 15, alltså
skilja på sexton register, och processorn lägger till 16: \code{0000} betyder \code{r16} och
\code{1111} betyder \code{r31}. Det finns inget bitmönster som betyder \code{r5}.

\pt{d}Ladda konstanten i ett övre register och kopiera den:

\begin{asmcode}
    ldi r16, 7                      ; r16 = 7
    mov r5, r16                     ; r5 = r16 = 7
\end{asmcode}

\noindent\code{mov} fungerar med alla 32 register, eftersom den inte har någon konstant att få
plats med.

\answer{c7:ex:2}{Tre minnen}
\pt{a}I programminnet, flashminnet.

\pt{b}I dataminnet, SRAM. Under SRAM, på dataminnets lägsta adresser, ligger också registren och
I/O-registren.

\pt{c}Flashminnet behåller innehållet utan ström. Därför börjar ett Arduino-kort köra sitt program
så fort det får ström, utan att programmet behöver laddas in igen. SRAM töms när strömmen bryts.

\pt{d}\code{RAMEND}, den sista adressen i SRAM. Där börjar stacken, som kapitel~\ref{c10:ch}
handlar om.

\answer{c7:ex:3}{Samma tal, olika skrivsätt}
\pt{a}45 = 32 + 8 + 4 + 1 = \code{0b00101101}, och \code{0010 1101} = \code{0x2D}:

\begin{asmcode}
    ldi r16, 45
    ldi r16, 0b00101101
    ldi r16, 0x2D
\end{asmcode}

\pt{b}\leavevmode

\begin{narrowtable}{@{}ccc@{}}
  \thead{Decimalt} & \thead{Binärt} & \thead{Hexadecimalt} \\
  \midrule
  255 & \code{0b11111111} & \code{0xFF} \\
  16 & \code{0b00010000} & \code{0x10} \\
  100 & \code{0b01100100} & \code{0x64} \\
\end{narrowtable}

\noindent 100 = 64 + 32 + 4, alltså bitarna 6, 5 och 2.

\pt{c}\leavevmode
\begin{itemize}
  \item \code{0b01010101} = 64 + 16 + 4 + 1 = \textbf{85}.
  \item \code{0x3F} = 3 · 16 + 15 = \textbf{63}.
  \item \code{$80} = 8 · 16 = \textbf{128}. \code{$} är ett annat sätt att skriva \code{0x}.
  \item \code{'a'} = \code{0x61} = 6 · 16 + 1 = \textbf{97}.
\end{itemize}

\pt{d}Binärt. Då står varje bit, och därmed varje lysdiod, för sig i programmet:
\code{0b10000001} tänder de två yttersta lysdioderna, vilket inte syns lika lätt i \code{0x81} och
inte alls i \code{129}.

\answer{c7:ex:4}{Förutsäg registren}
\begin{narrowtable}{@{}lccccc@{}}
  \thead{Efter} & \thead{\code{r16}} & \thead{\code{r17}} & \thead{\code{r18}} &
    \thead{\code{r19}} & \thead{\code{r20}} \\
  \midrule
  \code{ldi r16, 0x10} & 0x10 & 0x00 & 0x00 & 0x00 & 0x00 \\
  \code{ldi r17, 7} & 0x10 & 0x07 & 0x00 & 0x00 & 0x00 \\
  \code{mov r18, r17} & 0x10 & 0x07 & 0x07 & 0x00 & 0x00 \\
  \code{inc r18} & 0x10 & 0x07 & 0x08 & 0x00 & 0x00 \\
  \code{inc r18} & 0x10 & 0x07 & 0x09 & 0x00 & 0x00 \\
  \code{dec r16} & 0x0F & 0x07 & 0x09 & 0x00 & 0x00 \\
  \code{mov r19, r16} & 0x0F & 0x07 & 0x09 & 0x0F & 0x00 \\
  \code{clr r17} & 0x0F & 0x00 & 0x09 & 0x0F & 0x00 \\
  \code{ser r20} & 0x0F & 0x00 & 0x09 & 0x0F & 0xFF \\
  \code{dec r20} & 0x0F & 0x00 & 0x09 & 0x0F & 0xFE \\
\end{narrowtable}

\noindent Två rader brukar ge fel svar:
\begin{itemize}
  \item \textbf{\code{dec r16} från \code{0x10}} blir \code{0x0F}, inte \code{0x09}.
    Hexadecimalt räknar man ned från \code{10} till \code{0F}, precis som decimalt från 10 till 9.
    Den som svarar \code{0x09} har räknat hexadecimalt som om det vore decimalt.
  \item \textbf{\code{mov r19, r16}} kopierar det värde \code{r16} har \textbf{då}, \code{0x0F},
    inte det den fick i början.
\end{itemize}

\answer{c7:ex:5}{Hitta felen}
Sex fel:
\begin{enumerate}
  \item \textbf{\code{main}} saknar kolon. Utan kolon läser assemblern ordet som en instruktion,
    och det finns ingen instruktion som heter \code{main}.
  \item \textbf{\code{ldi r16 5}} saknar kommatecken mellan operanderna.
  \item \textbf{\code{ldi r10, 3}}: \code{ldi} når bara \code{r16}--\code{r31}. Använd till
    exempel \code{r20}.
  \item \textbf{\code{mov r17 r16}} saknar kommatecken.
  \item \textbf{\code{inc r20, 1}}: \code{inc} har bara en operand. Den ökar alltid med ett.
  \item \textbf{\code{rjmp mian}}: etiketten heter \code{main}. Assemblern känner inte till
    \code{mian}.
\end{enumerate}

\needspace{14\baselineskip}
Det rättade programmet:

\begin{asmcode}
.include "m328Pdef.inc"

.org 0x0000
    rjmp main

main:
    ldi r16, 5
    ldi r20, 3
    mov r17, r16
    inc r20
    rjmp main
\end{asmcode}

I Microchip Studio får du inte alltid alla sex felen på en gång. Rätta det första, bygg om, och
fortsätt med nästa, som i \secref{c7:c4}.

\answer{c7:ex:6}{Ditt första egna program}
\noanswer{Programmet självt har inget lösningsförslag här; simulatorn visar om det gör vad det
ska.}

\pt{a}\leavevmode

\begin{narrowtable}{@{}lcc@{}}
  \thead{Register} & \thead{Decimalt} & \thead{Hexadecimalt} \\
  \midrule
  \code{r16} & 12 & 0x0C \\
  \code{r17} & 15 & 0x0F \\
  \code{r18} & 0 & 0x00 \\
  \code{r19} & 170 & 0xAA \\
\end{narrowtable}

\pt{b}Värdena ska stämma. Om \code{r19} är 0 har \code{clr r18} hamnat före \code{mov r19, r18}:
ordningen spelar roll, eftersom \code{mov} kopierar det värde registret har just då.

\pt{c}Programmet har åtta instruktioner i \code{main}: \code{ldi}, \code{mov}, tre \code{inc},
\code{ldi}, \code{mov} och \code{clr}. \code{rjmp main} tar 2 cykler och de åtta instruktionerna 1
cykel var: 2 + 8 = \textbf{10} cykler. Cycle Counter ska visa 10 när den gula pilen står på
\code{rjmp end}.

\answer{c7:ex:7}{Maskinkod för hand}
\pt{a}Mönstret är \code{1110 KKKK dddd KKKK}, med konstantens höga halva i de första \code{K} och
den låga halvan i de sista.

\code{ldi r16, 0x2A}: konstanten \code{0x2A} = \code{0010 1010}, registret \code{r16} ger
\code{dddd} = 16 -- 16 = 0 = \code{0000}:

\begin{textdrawing}
1110  0010  0000  1010  =  0xE20A
\end{textdrawing}

\noindent De övriga på samma sätt:

\begin{narrowtable}{@{}lcccc@{}}
  \thead{Instruktion} & \thead{\code{K} höga halvan} & \thead{\code{dddd}} &
    \thead{\code{K} låga halvan} & \thead{Maskinkod} \\
  \midrule
  \code{ldi r16, 0x2A} & \code{0010} & \code{0000} (r16) & \code{1010} & \code{0xE20A} \\
  \code{ldi r24, 0x2A} & \code{0010} & \code{1000} (r24) & \code{1010} & \code{0xE28A} \\
  \code{ldi r31, 0xFF} & \code{1111} & \code{1111} (r31) & \code{1111} & \code{0xEFFF} \\
  \code{ldi r20, 0x03} & \code{0000} & \code{0100} (r20) & \code{0011} & \code{0xE043} \\
\end{narrowtable}

\noindent En genväg: i hexadecimal form är maskinkoden \code{E}, konstantens första siffra,
registrets nummer minus 16, och konstantens andra siffra. \code{ldi r24, 0x2A} blir
\code{E 2 8 A}.

\pt{b}Listfilen visar samma fyra ord, \code{e20a}, \code{e28a}, \code{efff} och \code{e043}, på
adresserna där instruktionerna hamnade.

\pt{c}De ska stämma. Det vanligaste felet är att skriva konstanten i ett stycke, till exempel
\code{1110 0000 0010 1010} = \code{0xE02A} för \code{ldi r16, 0x2A}. Det är fel ordning på fälten:
konstanten är delad, med registerfältet emellan. Ett annat vanligt fel är att skriva registrets
nummer utan att dra ifrån 16, \code{1 0000}, vilket inte ens får plats i fyra bitar.

\pt{d}Registerfältet har fyra bitar och betyder \code{r16}--\code{r31}. Det finns inget värde i
fältet som betyder \code{r15}, så det finns ingen maskinkod att skriva, och assemblern ger ett fel.

\answer{c7:ex:8}{Programmet som inte tog slut}
\pt{a}Processorn slutar aldrig köra. Det finns inget operativsystem att återvända till, så när
programmet har gjort sitt måste det stanna någonstans, och \code{rjmp end} hoppar till sig självt
för evigt.

\pt{b}Processorn fortsätter med det som ligger efter programmet i flashminnet. Där har du inte
skrivit något: ett tomt flashminne innehåller bara ettor, \code{0xFF}, och processorn försöker
tolka dem som instruktioner. I praktiken löper programräknaren då igenom resten av minnet, slår
runt till adress 0 och kör programmet från början igen, om och om igen. Programmet ser ut att
starta om av sig självt, vilket är ett förvirrande fel att leta efter om man inte vet om det.

\pt{c}Adresserna närmast 0 är reserverade för avbrottsvektorer. Boken använder inga avbrott, men
vanan att börja med ett hopp gör att programmet fortsätter att fungera om något läggs där, och det
är så nästan alla AVR-program börjar (\secref{c7:b6}).

\answer{c7:ex:9}{Tillåtet eller inte?}
\pt{a}\textbf{Tillåtet.} \code{mov} kopierar mellan två valfria register, \code{r0}--\code{r31}.

\pt{b}\textbf{Inte tillåtet.} \code{ldi} når bara \code{r16}--\code{r31}. Skriv
\code{ldi r16, 16} och \code{mov r5, r16}.

\pt{c}\textbf{Inte tillåtet.} \code{mov} kopierar mellan register; den har ingen konstant. Skriv
\code{ldi r16, 5}.

\pt{d}\textbf{Inte tillåtet.} \code{ldi} laddar en konstant, inte ett register. Skriv
\code{mov r16, r5}.

\pt{e}\textbf{Tillåtet.} \code{inc} fungerar med alla 32 register.

\answer{c7:ex:10}{Hur lång tid tar programmet?}
\pt{a}\code{rjmp main} tar 2 cykler och de tio instruktionerna i \code{main} 1 cykel var:
2 + 10 = \textbf{12} cykler.

\pt{b}12 · 62,5~ns = \textbf{750~ns}, det vill säga 0,75 mikrosekunder.

\pt{c}1~s / 750~ns $\approx$ 1,3 miljoner gånger per sekund. Så snabbt går det; det är därför
programmet måste stegas i simulatorn för att man ska se något.

\answer{c7:ex:11}{Maskinkod baklänges}
Dela ordet i fyra hexadecimala siffror: \code{E}, konstantens höga halva, registrets nummer minus
16, och konstantens låga halva.

\pt{a}\code{0xE3F5}: konstanten \code{0x35}, registerfältet \code{F} = 15, 15 + 16 = 31.
\textbf{\code{ldi r31, 0x35}}.

\pt{b}\code{0xE0A0}: konstanten \code{0x00}, registerfältet \code{A} = 10, 10 + 16 = 26.
\textbf{\code{ldi r26, 0x00}}.

\pt{c}\code{0xEF0F}: konstanten \code{0xFF}, registerfältet \code{0} = 0, alltså \code{r16}.
\textbf{\code{ldi r16, 0xFF}}. Det är exakt vad \code{ser r16} gör, och \code{ser} är i själva
verket bara ett annat namn på \code{ldi} med konstanten \code{0xFF}: assemblern gör samma maskinkod
av båda. Det syns i listfilen i \secref{c7:b7}, där \code{ser r20} blev \code{ef4f}.
